\scanpage{39}

le

\uline{Corollaire} Supposons que les $Z_i$ soient \uline{\unsure{rares}} dans $S$. Alors il existe un $x\in S$, avec $\dim\overline{\{x\}}=1$, tel que $x\notin\bigcup Z_i$.

Remplaçant $S$ par $T=\operatorname{Spec} V[[T_1,\ldots,T_{\unsure{r|n}}]]$, \deleted{\unsure{on a}} on voit que les $T_i:=\operatorname{Im}(Z_{\unsure{i}}\unsure{-|\to}T)$ sont rares, et il suffit de trouver un $x\in T$, avec $\dim\overline{\{x\}}=1$, tel que $x\notin\bigcup T_i$. (En effet, il existe alors un $y\in S_{\unsure{2|x}}$ au dessus de $x$, et on a nécessairement $\dim\overline{\{y\}}=1$).

On va montrer qu'il existe $\unsure{\overline{\mathrm{m}}|\overline{\mathrm{me}}}$ un $x\in T$ provenant d'un $V$-hom
\[
V[[T_1,\ldots,T_r]]\unsure{-|\longrightarrow}V\;.\qquad \deleted{\text{\unsure{Comme}}}\unsure{\text{,}}
\]

En effet, l'ens. des ces hom s'identifie à $\mathfrak m^{\unsure{r|n}}$ ($\mathfrak m=\operatorname{rad}(\unsure{V|W})$). Il a une structure d'espace métrique complet. Montrons que pour tte partie fermée rare $Z$ de $T$, l'ens des $\varphi$ tels que $V(\operatorname{Ker}\varphi)\subset Z$ est rare. En effet, on peut supposer que \deleted{$Z$ est \unsure{défi|défini}} $Z=V(f)$, \inserted{$f\in V[[T_1\unsure{,|-}\ldots T_r]]$, $f\neq0$,} et $\lvert V(\operatorname{Ker}\varphi)\rvert\subset\lvert V(f)\rvert$ équivaut (comme $V(\operatorname{Ker}\varphi)$ est intègre \unsure{bien|donc} évidemment) à $V(\operatorname{Ker}\varphi)\subset V(f)$ \unsure{i.e.|\ldots} $f\in\operatorname{Ker}\varphi$ i.e. $\varphi(f)=0$; i.e. $f(\varphi_1,\ldots,\varphi_r)=0$. Il est immédiat que c'est une partie fermée rare.
% archive sheet number: 38
